\section{Limit transition to noncolliding walks}\label{sec:limit_transition}

In this section we perform the limit transition as $N\to+\infty$ in the
Markov kernels $\Lambda^{N}_{N-1}$ \eqref{eq:traditional_links_Lambda}
under the principal specialization
\eqref{eq:principal_spec}, and prove
Theorem~\ref{thm:Macdonald_noncolliding_intro}.

\subsection{Setup}\label{sub:limit_transition_setup}

Denote by $\mathbb{W}_m$ the space of $m$-particle configurations
in $\mathbb{Z}_{\ge0}$, that is,
\begin{equation}
	\label{eq:W_m_space}
	\mathbb{W}_m\coloneqq
	\{
		\vec{\mathsf{x}}=(\mathsf{x}_1>\mathsf{x}_2>\dots>\mathsf{x}_m\ge0 )
	\}\subset \mathbb{Z}_{\ge0}^{m}.
\end{equation}
By $\mathbb{W}_m(N)$ denote the finite subset of $\mathbb{W}_m$ determined by
the condition
$\mathsf{x}_1\le N+m-2$. Define the injective maps
\begin{equation}
	\label{eq:pi_maps_1}
	\pi\colon \ \mathbb{W}_m(N)\to \mathbb{Y}(N), \qquad
	\overline\pi\colon \ \mathbb{W}_m(N)\to \mathbb{Y}(N-1),
\end{equation}
as follows. If $\lambda=\pi(\vec{\mathsf{x}})\in \mathbb{Y}(N)$
and $\mu=\overline\pi(\vec {\mathsf{y}})\in \mathbb{Y}(N-1)$, then
\begin{gather}
 \{ \lambda_1-1,\lambda_2-2,\dots,\lambda_N-N \}=
 \{ 0,1,2,\dots,N+m-1 \}\setminus
		\{\mathsf{x}_1,\dots, \mathsf{x}_m \},\nonumber\\
		 \{ \mu_1-1,\mu_2-2,\dots,\mu_{N-1}-(N-1) \}=
		 \{ 1,2,\dots,N+m-1 \}\setminus
		\{\mathsf{y}_1+1,\dots, \mathsf{y}_m+1 \}.\!\!\!\label{eq:pi_maps_2}
	\end{gather}
For fixed $\vec{\mathsf{x}}$ and growing $N$,
almost all parts of $\lambda=\pi(\vec{\mathsf{x}})$
are equal to $N+m$, and there is a defect in a few last parts of
$\lambda$. The columns of this defect are
encoded through
$\vec{\mathsf{x}}$. A similar description holds for
$\mu=\overline\pi(\vec{\mathsf{y}})$.
In multiplicative notation for partitions, we have
\begin{gather}
		\lambda=\pi(\vec{\mathsf{x}})=
		(N+m)^{N-\mathsf{x}_1+m-1}(N+m-1)^{\mathsf{x}_1-\mathsf{x}_2-1}
		\cdots (N+1)^{\mathsf{x}_{m-1}-\mathsf{x}_m-1} N^{\mathsf{x}_m},
\nonumber\\
		\mu=\overline\pi(\vec{\mathsf{y}})=
		(N+m)^{N-\mathsf{y}_1+m-2}(N+m-1)^{\mathsf{y}_1-\mathsf{y}_2-1}
		\cdots (N+1)^{\mathsf{y}_{m-1}-\mathsf{y}_m-1}
		N^{\mathsf{y}_m}.\label{eq:pi_maps_multiplicative_notation}
	\end{gather}
See
Figure~\ref{fig:pi_of_x} for an illustration.

\begin{figure}[htpb]
	\centering
	\includegraphics[width=.8\textwidth]{fig_pi_of_x.pdf}
	\caption{Left: Young diagram of $\pi(\vec{\mathsf{x}})$ defined by
	\eqref{eq:pi_maps_1}--\eqref{eq:pi_maps_2}.
	The column lengths of the defect are $\mathsf{x}_j-(m-j)$.
	Right: Young diagram of $\overline\pi(\vec{\mathsf{y}})$.
	Observe that
	$\lambda=\pi(\vec{\mathsf{x}})$ and
	$\overline\pi(\vec{\mathsf{y}})$ differ only by
	adding the first row.}	\label{fig:pi_of_x}
\end{figure}

\begin{Lemma}	\label{lemma:interlacing}
	For any $\vec{\mathsf{x}},\vec{\mathsf{y}}\in \mathbb{W}_m(N)$,
	we have $\overline\pi(\vec{\mathsf{y}})\prec \pi(\vec{\mathsf{x}})$
	if and only if $\mathsf{y}_i=\mathsf{x}_i$ or $\mathsf{y}_i=\mathsf{x}_i-1$
	for all $i=1,\dots,m $.
\end{Lemma}
\begin{proof}
	Straightforward verification.
\end{proof}

In the rest of this section we
fix arbitrary $\vec{\mathsf{x}},\vec{\mathsf{y}}\in \mathbb{W}_m$ and
compute the limit of
$\Lambda^{N}_{N-1}(\pi(\vec{\mathsf{x}}),\overline\pi(\vec{\mathsf{y}}))$
\eqref{eq:traditional_links_Lambda}
under principal
specialization $x_j=t^{j-1}$
as $N\to+\infty$. Clearly, for sufficiently large $N$ the
partitions
$\pi(\vec{\mathsf{x}}),\overline\pi(\vec{\mathsf{y}})$
are well-defined.
We will show that this limit is equal
to the Markov kernel $\Upsilon_m(\vec{\mathsf{x}},\vec{\mathsf{y}})$
defined by~\eqref{eq:intro_Upsilon}.

Throughout the computation we adopt the convention
that
$\lambda=\pi(\vec{\mathsf{x}})$, $\mu=\overline\pi(\vec{\mathsf{y}})$.

\subsection{Initial Markov kernel}\label{sub:computation_initial_formula}

Our starting point is the formula for
$\Lambda^{N}_{N-1}(\pi(\vec{\mathsf{x}}),\overline\pi(\vec{\mathsf{y}}))$, see
\eqref{eq:traditional_links_Lambda},
under the principal
specialization $x_j=t^{j-1}$, which takes the form (where $\mu\prec\lambda$):
\begin{gather}
	\nonumber
		\Lambda^{N}_{N-1}(\lambda,\mu)
		=
		\psi_{\lambda/\mu}(q,t)\hspace{1pt} x_N^{|\lambda|-|\mu|}\hspace{1pt}
		\frac{
		P_\mu(x_1,\dots,x_{N-1} \mid q,t )
		}
		{P_\lambda(x_1,\dots,x_N \mid q,t)}
		\\
		\nonumber
\hphantom{\Lambda^{N}_{N-1}(\lambda,\mu)}{}
=
		\psi_{\lambda/\mu}
		(q,t)\hspace{1pt}
		t^{(N-1)(|\lambda|-|\mu|)}\hspace{1pt}
		\frac{P_\mu\big(1,t,\dots,t^{N-2}\mid q,t \big)}{P_\lambda\big(1,t,\dots,t^{N-2},t^{N-1}\mid q,t\big)}
		\\
		\nonumber
\hphantom{\Lambda^{N}_{N-1}(\lambda,\mu)}{}
=
		t^{(N-1)(|\lambda|-|\mu|)}
		\prod_{1\le i<j\le N-1}
		\frac{\mathsf{f}\big(q^{\mu_i-\mu_j}t^{j-i}\big)\hspace{1pt}
		\mathsf{f}(q^{\lambda_i-\lambda_{j+1}}t^{j-i})}
		{\mathsf{f}\big(q^{\lambda_i-\mu_j}t^{j-i}\big)\hspace{1pt}
		\mathsf{f}\big(q^{\mu_i-\lambda_{j+1}}t^{j-i}\big)}
		\\
\hphantom{\Lambda^{N}_{N-1}(\lambda,\mu)=}{}
\times
		t^{n(\mu)}\prod_{\square\in \mu}
		\frac{1-q^{a'(\square)}t^{N-1-l'(\square)}}
		{1-q^{a(\square)}t^{l(\square)+1}}
		t^{-n(\lambda)}
		\prod_{\square\in \lambda}
		\frac
		{1-q^{a(\square)}t^{l(\square)+1}}
		{1-q^{a'(\square)}t^{N-l'(\square)}}.		\label{eq:initial_formula}
\end{gather}
Here we used
\eqref{eq:principal_spec} and
\eqref{eq:psi_Pieri}.
Our next steps are devoted to
taking the limit as
$N\to+\infty$
in various parts of the product~\eqref{eq:initial_formula}.

\subsection[Power of t]{Power of $\boldsymbol{t}$}

Let us first consider the overall power of $t$
in \eqref{eq:initial_formula}
which is equal to
$t^{(N-1)(|\lambda|-|\mu|)+n(\mu)-n(\lambda)}$.
Our aim is to express quantities
depending on $\lambda$, $\mu$ through
$\vec{\mathsf{x}}$, $\vec{\mathsf{y}}$.
Adopt the convention
$\mathsf{x}_0=N+m$,
$\mathsf{y}_0=N+m-1$,
$\mathsf{x}_{m+1}=\mathsf{y}_{m+1}=-1$,
so that
\begin{equation*}
	N+m-\mathsf{x}_1-1=\mathsf{x}_0-\mathsf{x}_1-1,\qquad
	N+m-\mathsf{y}_1-2=\mathsf{y}_0-\mathsf{y}_1-1.
\end{equation*}
We have from \eqref{eq:pi_maps_multiplicative_notation}:
\begin{gather*}
	|\lambda|-|\mu|	=
	\sum_{i=0}^{m}
	(N+i)
	(
	( \mathsf{x}_{m-i}-\mathsf{x}_{m-i+1}-1 )-
	( \mathsf{y}_{m-i}-\mathsf{y}_{m-i+1}-1)
)
	\\
\hphantom{|\lambda|-|\mu|}{}
=
	\sum_{i=0}^{m}
	(N+i)
	(
	( \mathsf{x}_{m-i}-\mathsf{x}_{m-i+1} )-
	( \mathsf{y}_{m-i}-\mathsf{y}_{m-i+1} )
)
	\\
\hphantom{|\lambda|-|\mu|}{}
=
	N\sum_{i=0}^{m}
	(
	( \mathsf{x}_{m-i}-\mathsf{x}_{m-i+1} )-
	( \mathsf{y}_{m-i}-\mathsf{y}_{m-i+1} )
	)\\
\hphantom{|\lambda|-|\mu|=}{}
	+
	\sum_{i=0}^{m}i
	(
	( \mathsf{x}_{m-i}-\mathsf{x}_{m-i+1} )-
	( \mathsf{y}_{m-i}-\mathsf{y}_{m-i+1})
	)
	\\
\hphantom{|\lambda|-|\mu|}{}
=
	N+m+
	|\vec {\mathsf{y}}|-|\vec{\mathsf{x}}|.
\end{gather*}
Moreover, from
\eqref{eq:n_lambda_notation} and \eqref{eq:pi_maps_multiplicative_notation}
we have
\begin{gather}
	\nonumber
	n(\mu)-n(\lambda)
	=
	\sum_{j=0}^{m}
	(N+m-j)
	\left[
		\binom{\mathsf{y}_0-\mathsf{y}_{j+1}-(j+1)}2-
		\binom{\mathsf{y}_0-\mathsf{y}_{j}-j}2
	\right]
	\\
\hphantom{n(\mu)-n(\lambda)=}{}
-
	\sum_{j=0}^{m}
	(N+m-j)
	\left[
		\binom{\mathsf{x}_0-\mathsf{x}_{j+1}-(j+1)}2-
		\binom{\mathsf{x}_0-\mathsf{x}_{j}-j}2
	\right]
\nonumber\\
\hphantom{n(\mu)-n(\lambda)}{}
=
	-(N-1)(N+m)
	-\binom{m}{2}
	+|\vec{\mathsf{x}}|\nonumber\\
\hphantom{n(\mu)-n(\lambda)=}{}
+
	\frac12
	\sum_{i=1}^{m}
	(\mathsf{y}_i-\mathsf{x}_i)(\mathsf{x}_i+\mathsf{y}_i+3+2i-2m-2N).
\label{eq:newlabel}
\end{gather}
Indeed, the coefficients
by $\mathsf{x}_i$, $\mathsf{x}_i^2$, $\mathsf{y}_i$, $\mathsf{y}_i^2$, $1\le i\le m$, in the right-hand side of~\eqref{eq:newlabel}
are, respectively,
\begin{equation*}
	N - i + m - \frac{1}{2}
	,\
	-\frac{1}{2}
	,\
	-N + i - m + \frac32
	,\
	\frac{1}{2},
\end{equation*}
which are the same as in the left-hand side.
The free term in the left-hand side is
$-\binom m2+m-N^2-(m-2)N$,
which is readily matched to the free term in the
right-hand side by virtue of
our conventions about
$\mathsf{x}_0$, $\mathsf{y}_0$, $\mathsf{x}_{m+1}$, $\mathsf{y}_{m+1}$.

Let us further simplify the left-hand side of
\eqref{eq:newlabel}.
The $i$-th term in this sum
is rewritten as
\begin{gather*}
	\frac12(\mathsf{y}_i-\mathsf{x}_i)(\mathsf{x}_i+\mathsf{y}_i+3+2i-2m-2N)\\
\qquad{}
	=
	-(N-1)(\mathsf{y}_i-\mathsf{x}_i)
	+
	\frac12(\mathsf{y}_i-\mathsf{x}_i)(-\mathsf{x}_i+\mathsf{y}_i+1)
	+
	(\mathsf{y}_i-\mathsf{x}_i)(\mathsf{x}_i-m+i).
\end{gather*}
Since $\mathsf{y}_i-\mathsf{x}_i$ is equal to $0$ or $-1$,
the quantity
$(\mathsf{y}_i-\mathsf{x}_i)(-\mathsf{x}_i+\mathsf{y}_i+1)$
is identically zero.
Thus, we have
\begin{align*}
	n(\mu)-n(\lambda)=
	-(N-1)(N+m+|\vec{\mathsf{y}}|-|\vec{\mathsf{x}}|)
	+|\vec{\mathsf{x}}|-\binom m2
	+\sum_{i=1}^{m}
	(\mathsf{x}_{i}-m+i)(\mathsf{y}_{i}-\mathsf{x}_{i}).
\end{align*}
We see that
the $N$-dependent terms
in
$(N-1)(|\lambda|-|\mu|)+n(\mu)-n(\lambda)$
cancel out, and the overall
factor containing the
power of $t$
in
$\Lambda^{N}_{N-1}(\pi(\vec{\mathsf{x}}),\overline\pi(\vec{\mathsf{y}}))$
has the form
\begin{equation*}%\label{eq:power_of_t_final}
	t^{-\binom m2+|\vec{\mathsf{x}}|
	+\sum_{i=1}^{m}
	(\mathsf{x}_{i}-m+i)(\mathsf{y}_{i}-\mathsf{x}_{i})}.
\end{equation*}

\subsection{Coarms and colegs}
Addressing the factors in
\eqref{eq:initial_formula}
containing coarms and colegs of $\lambda$ and $\mu$,
we obtain using \eqref{eq:arms_legs}:
\begin{equation*}
	\prod_{\square\in \lambda}\big(1-q^{a'(\square)}t^{N-l'(\square)}\big)
	=
	\prod_{i=1}^{N}
	\prod_{j=1}^{\lambda_i}
	\big(1-q^{j-1}t^{N-i+1}\big)
	=
	\prod_{i=1}^{N}
	\big(t^{N+1-i};q\big)_{\lambda_i}.
\end{equation*}
This product is in the denominator, and a similar
factor
$\prod_{i=1}^{N-1}
\big(t^{N-i};q\big)_{\mu_i}$
appears in the numerator.
Let us show that the contribution coming from
coarms and colegs goes to one, that is,%
\begin{equation}
	\label{eq:coarm_coleg_result}
	\lim_{N\to+\infty}
	\frac{\prod_{i=1}^{N-1}
	\big(t^{N-i};q\big)_{\mu_i}}
	{\prod_{i=1}^{N}
	\big(t^{N+1-i};q\big)_{\lambda_i}}=1.
\end{equation}
To see this, observe that
$\mu_{i}-\lambda_{i+1}$ for all $i=1,\dots,N-1 $ is a nonnegative integer
which does not grow with $N$ as long as $N-i$ is fixed,
see~\eqref{eq:pi_maps_multiplicative_notation}. Therefore,
\begin{gather*}
	\frac{\big(t^{N-i};q\big)_{\mu_i}}{\big(t^{N+1-(i+1)};q\big)_{\lambda_{i+1}}}=
	\frac{\big(t^{N-i};q\big)_{\mu_i}}{\big(t^{N-i};q\big)_{\lambda_{i+1}}}\\
\hphantom{\frac{\big(t^{N-i};q\big)_{\mu_i}}{\big(t^{N+1-(i+1)};q\big)_{\lambda_{i+1}}}}{}
=
	\big(1-t^{N-i}q^{\lambda_{i+1}}\big)\big(1-t^{N-i}q^{\lambda_{i+1}+1}\big)\cdots\big(1-t^{N-i}q^{\mu_i-1}\big) ,
\end{gather*}
where the product in the right-hand side is finite. Since $\lambda_{i+1}$ and $\mu_i$ go to infinity
as $N\to\infty$, this product converges to~$1$.
There is one more factor in the denominator of~\eqref{eq:coarm_coleg_result},
namely,
\begin{equation*}
	\frac{1}{\big(t^N;q\big)_{\lambda_1}}=
	\frac{\big(t^Nq^{N+m};q\big)_\infty}{\big(t^N;q\big)_\infty}.
\end{equation*}
This factor also goes to $1$ as $N\to+\infty$,
and so the limit~\eqref{eq:coarm_coleg_result} is established.

\subsection{Arms and legs}

We now consider the product
in \eqref{eq:initial_formula}
involving arms and legs:
\begin{equation}
	\label{eq:arm_leg_product_initial}
	\prod_{\square\in \mu}
	\frac{1}
	{1-q^{a(\square)}t^{l(\square)+1}}
	\prod_{\square\in \lambda}
	\bigl(1-q^{a(\square)}t^{l(\square)+1}\bigr).
\end{equation}
Recall the notation
$\lambda=\pi(\vec{\mathsf{x}})$,
$\mu=\overline\pi(\vec{\mathsf{y}})$, see~\eqref{eq:pi_maps_multiplicative_notation}.
Observe that the Young diagrams
$\lambda=\pi(\vec{\mathsf{x}})$ and
$\overline\pi(\vec{\mathsf{y}})$ differ only by
adding the first row (cf.\ Figure~\ref{fig:pi_of_x}).
For each box $\square$ in the first row of $\lambda$, the quantity
$l(\square)$ is of order $N$, and so
$1-q^{a(\square)}t^{l(\square)+1}$ is close to~$1$
for large $N$.
Therefore, the product
\eqref{eq:arm_leg_product_initial}
has the same limit as
\begin{equation}	\label{eq:arm_leg_product_proof_1}
	\prod_{\square\in \overline\pi(\vec{\mathsf{y}})}
	\frac{1}
	{1-q^{a(\square)}t^{l(\square)+1}}
	\prod_{\square\in \overline\pi(\vec{\mathsf{x}})}
	\bigl(1-q^{a(\square)}t^{l(\square)+1}\bigr).
\end{equation}

\begin{Proposition}	\label{prop:arm_leg_main_result}
	As $N\to+\infty$, the product \eqref{eq:arm_leg_product_proof_1}
	converges to
	\begin{equation*}%\label{eq:arm_leg_limit_thing}
		\frac{V_{t,q}(\vec{\mathsf{y}})}{V_{t,q}(\vec{\mathsf{x}})}
		\prod_{i=1}^{m}
		\big( 1-q^{m-i}t^{\mathsf{x}_i-m+i}\mathbf{1}_{\mathsf{y}_i=\mathsf{x}_i-1} \big),
	\end{equation*}
	where $V_{t,q}$ is the $(t,q)$-Vandermonde given by \eqref{eq:V_tq}.
\end{Proposition}
\begin{proof}
	Recall the notation
	for the
	$(q,t)$-deformed hook polynomials
	\cite[formulas~(VI.8.1) and (VI.8.1$'$)]{Macdonald1995}:
	\begin{equation*}
		c_\nu(q,t)=
		\prod_{\square\in \nu}\bigl(1-q^{a(\square)}t^{l(\square)+1}\bigr),\qquad
		c'_\nu(q,t)=
		\prod_{\square\in \nu}\bigl(1-q^{a(\square)+1}t^{l(\square)}\bigr),
	\end{equation*}
	and observe that $c_\nu(q,t)=c'_{\nu'}(t,q)$, where $\nu'$ is the
	transposed Young diagram.
	Note that, in multiplicative notation,
	\begin{gather*}
		\overline\pi(\vec{\mathsf{x}})'=
		(N-1)^{N}
		(N-1-\mathsf{x}_m)
		(N-1-\mathsf{x}_{m-1}+1)
		(N-1-\mathsf{x}_{m-2}+2)
		\cdots\\
\hphantom{\overline\pi(\vec{\mathsf{x}})'=}{}
\times (N-1-\mathsf{x}_{1}+m-1)
	\end{gather*}
	and similarly for $\overline\pi(\vec{\mathsf{y}})'$,
	see Figure~\ref{fig:pi_of_x}.
	Adopt the convention
	$\mathsf{x}_{m+j}=\mathsf{y}_{m+j}
	=-j$ for $j=1,2,\dots $.
	Then we may shift the indices $j$ to encode the string $\overline\pi(\vec{\mathsf{x}})'$
	as
	$\overline\pi(\vec{\mathsf{x}})'_j=N+j-1-\mathsf{x}_{m-j}$,
	where $-N\le j\le m-1$, and similarly for $\vec{\mathsf{y}}$.
	The arm and leg lengths do not change under this shifting.
	This shift allows to directly refer to a known identity, the first one in
	\cite[Proposition~3.2]{kaneko1996q}, and write
	\begin{gather*}
		c_{\overline\pi(\vec{\mathsf{x}})}(q,t)
		=
		c'_{\overline\pi(\vec{\mathsf{x}})'}(t,q)
		\\
\hphantom{c_{\overline\pi(\vec{\mathsf{x}})}(q,t)}{}
=
		\frac{(t;t)_{\infty}^{N+m}}
		{
			\prod_{j=-N}^{m-1}
			\big(t^{N+j-\mathsf{x}_{m-j}}q^{m-1-j};t\big)_{\infty}
		}
		\prod_{-N\le i<j\le m-1}
		\frac{\big(q^{j-i}t^{i-j-\mathsf{x}_{m-i}+\mathsf{x}_{m-j}+1};t\big)_{\infty}}
		{\big(q^{j-i-1}t^{i-j-\mathsf{x}_{m-i}+\mathsf{x}_{m-j}+1};t\big)_{\infty}}
		\\
\hphantom{c_{\overline\pi(\vec{\mathsf{x}})}(q,t)}{}
		=
		\frac{(t;t)_{\infty}^{N+m}}
		{
			\prod_{j=-N}^{m-1}
			\big(t^{N+j-\mathsf{x}_{m-j}}q^{m-1-j};t\big)_{\infty}
		}
		\prod_{-N\le i<j\le -1}
		\frac{\big(q^{j-i}t;t\big)_{\infty}}
		{\big(q^{j-i-1}t;t\big)_{\infty}}
		\\
\hphantom{c_{\overline\pi(\vec{\mathsf{x}})}(q,t)=}{}
\times
		\prod_{\substack{-N\le i\le -1 \\ 0\le j\le m-1}}
		\frac{\big(q^{j-i}t^{-j+\mathsf{x}_{m-j}+1};t\big)_{\infty}}
		{\big(q^{j-i-1}t^{-j+\mathsf{x}_{m-j}+1};t\big)_{\infty}}
		\underbrace{
		\prod_{1\le i<j\le m}
		\frac{\big(q^{j-i}t^{i-j-\mathsf{x}_{j}+\mathsf{x}_{i}+1};t\big)_{\infty}}
		{\big(q^{j-i-1}t^{i-j-\mathsf{x}_{j}+\mathsf{x}_{i}+1};t\big)_{\infty}}}
		_{1/V_{t,q}(\vec{\mathsf{x}})}.
	\end{gather*}
	Therefore,
	the product \eqref{eq:arm_leg_product_proof_1}
	becomes
	\begin{gather*}
		\frac{c_{\overline\pi(\vec{\mathsf{x}})}(q,t)}
		{c_{\overline\pi(\vec{\mathsf{y}})}(q,t)}
		=
		\frac{V_{t,q}(\vec{\mathsf{y}})}
		{V_{t,q}(\vec{\mathsf{x}})}
			\prod_{j=1}^{m}
		\frac{
			\big(t^{N+m-j-\mathsf{y}_{j}}q^{j-1};t\big)_{\infty}
		}
		{
			\big(t^{N+m-j-\mathsf{x}_{j}}q^{j-1};t\big)_{\infty}
		}\\
\hphantom{\frac{c_{\overline\pi(\vec{\mathsf{x}})}(q,t)}
		{c_{\overline\pi(\vec{\mathsf{y}})}(q,t)}=}{}
\times
		\prod_{j=1}^{m}
		\prod_{i=-N}^{-1}
		\frac{\big(q^{m-j-i}t^{j-m+\mathsf{x}_{j}+1};t\big)_{\infty}}
		{\big(q^{m-j-i-1}t^{j-m+\mathsf{x}_{j}+1};t\big)_{\infty}}
		\frac
		{\big(q^{m-j-i-1}t^{j-m+\mathsf{y}_{j}+1};t\big)_{\infty}}
		{\big(q^{m-j-i}t^{j-m+\mathsf{y}_{j}+1};t\big)_{\infty}}.
	\end{gather*}
	The first product over $1\le j\le m$
	converges to $1$ as $N\to+\infty$
	thanks to the presence of $t^N$.
	In the second product over $1\le j\le m$, the terms
	where $\mathsf{y}_j=\mathsf{x}_j$ are simply equal to $1$.
	When $\mathsf{y}_j=\mathsf{x}_j-1$, we have
	\begin{gather*}
		\prod_{i=-N}^{-1}
		\frac{\big(q^{m-j-i}t^{j-m+\mathsf{x}_{j}+1};t\big)_{\infty}}
		{\big(q^{m-j-i-1}t^{j-m+\mathsf{x}_{j}+1};t\big)_{\infty}}
		\frac
		{\big(q^{m-j-i-1}t^{j-m+\mathsf{x}_{j}};t\big)_{\infty}}
		{\big(q^{m-j-i}t^{j-m+\mathsf{x}_{j}};t\big)_{\infty}}
	\\
\qquad{}	 =
		\prod_{i=1}^{N}
		\frac{1-q^{m-j+i-1}t^{j-m+\mathsf{x}_{j}}}
		{1-q^{m-j+i}t^{j-m+\mathsf{x}_{j}}}
	 =
		\frac{1-q^{m-j}t^{j-m+\mathsf{x}_{j}}}
		{1-q^{m-j+N}t^{j-m+\mathsf{x}_{j}}},
	\end{gather*}
	and the $N\to+\infty$ limit of this expression is
	$1-q^{m-j}t^{\mathsf{x}_{j}-m+j}$.
	This completes the proof.
\end{proof}

\subsection{Pieri coefficient}

It remains to consider the $N\to+\infty$ limit of the Pieri
coefficient $\psi_{\lambda/\mu}(q,t)$
entering \eqref{eq:initial_formula}.
Recall the convention
$\mathsf{x}_{m+j}=\mathsf{y}_{m+j}
=-j$ for $j=1,2,\dots $,
and encode (shifting the indices)
\begin{equation*}
	\lambda'_i=\pi(\vec{\mathsf{x}})'_i=N+i-\mathsf{x}_{m-i},
	\qquad
	\mu'_j=\overline\pi(\vec{\mathsf{y}})'_j=N-1+j-\mathsf{y}_{m-j}.
\end{equation*}
We can now use
the well-known duality
\cite[proof of formula~(6.24)]{Macdonald1995}
\begin{equation*}
	\psi_{\lambda/\mu}(q,t)=
	\psi'_{\lambda'/\mu'}(t,q),
\end{equation*}
and obtain from~\eqref{eq:psi_prime_Pieri}:
\begin{gather*}
	\psi'_{\lambda'/\mu'}(t,q)
	=
	\prod_{\substack{-N\le i<j\le m-1
	\\
	\lambda'_{i}=\mu'_{i},\, \lambda'_{j}=\mu'_{j}+1}}
	\frac{\big(1-t^{\mu'_{i}-\mu'_{j}} q^{j-i-1}\big)\big(1-t^{\lambda'_{i}-\lambda'_{j}} q^{j-i+1}\big)}
	{\big(1-t^{\mu'_{i}-\mu'_{j}} q^{j-i}\big)\big(1-t^{\lambda'_{i}-\lambda'_{j}} q^{j-i}\big)}
	\\
\hphantom{\psi'_{\lambda'/\mu'}(t,q)=}{}
=
	\prod_{\substack{-N\le i<j\le m-1
	\\
	\mathsf{y}_{m-i}=\mathsf{x}_{m-i}-1,\\
	\mathsf{y}_{m-j}=\mathsf{x}_{m-j}}}
	\frac{\big(1-t^{i-j+\mathsf{y}_{m-j}-\mathsf{y}_{m-i}}
	q^{j-i-1}\big)
	\big(1-t^{i-j+\mathsf{x}_{m-j}-\mathsf{x}_{m-i}} q^{j-i+1}\big)}
	{\big(1-t^{i-j+\mathsf{y}_{m-j}-\mathsf{y}_{m-i}} q^{j-i}\big)
	\big(1-t^{i-j+\mathsf{x}_{m-j}-\mathsf{x}_{m-i}} q^{j-i}\big)}.
\end{gather*}
Split the product over $i<j$ into three parts.
The first part with $-N\le i<j\le -1$
cancels out and is equal to $1$
since
we have $\mathsf{y}_{m-i}=\mathsf{x}_{m-i}$ for all $-N\le i\le -1$.
The second part
with $-N\le i\le -1$ and $0\le j\le m-1$ is also equal to $1$ for the same reason.
The third part with $0\le i<j\le m-1$ equals
\begin{equation}
	\label{eq:psi_prime_in_limit}
	\prod_{\substack{1\le i<j\le m
	\\
	\mathsf{y}_{i}=\mathsf{x}_{i}
	,\,
	\mathsf{y}_{j}=\mathsf{x}_{j}-1}}
	\frac{\big(1-t^{i-j+\mathsf{y}_{i}-\mathsf{y}_{j}}
	q^{j-i-1}\big)
	\big(1-t^{i-j+\mathsf{x}_{i}-\mathsf{x}_{j}} q^{j-i+1}\big)}
	{\big(1-t^{i-j+\mathsf{y}_{i}-\mathsf{y}_{j}} q^{j-i}\big)
	\big(1-t^{i-j+\mathsf{x}_{i}-\mathsf{x}_{j}} q^{j-i}\big)},
\end{equation}
which is independent of $N$. Note that
\eqref{eq:psi_prime_in_limit} is equal to a Pieri coefficient
$\psi'_{(\vec{\mathsf{x}}-\delta_m)/(\vec{\mathsf{y}}-\delta_m)}(t,q)$,
where $\delta_m=(m-1,m-2,\dots,2,1,0)$ is the
staircase partition.

\subsection{Final result}

Putting together all the computations from this section, we see that the
$N\to+\infty$ limit of the Markov kernel $\Lambda^{N}_{N-1}(\pi(\vec{\mathsf{x}}),
\overline\pi(\vec{\mathsf{y}}))$ is the Markov kernel (on $\mathbb{W}_m$)
of the Macdonald noncolliding walks with the absorbing wall at zero:
\begin{gather}
	\nonumber
	\Upsilon_m(\vec{\mathsf{x}},\vec{\mathsf{y}}) =
	t^{-\binom m2+|\vec{\mathsf{x}}|
	+\sum_{i=1}^{m}
	(\mathsf{x}_{i}-m+i)(\mathsf{y}_{i}-\mathsf{x}_{i})}	
	\hspace{1pt}
	\frac{V_{t,q}(\vec{\mathsf{y}})}{V_{t,q}(\vec{\mathsf{x}})}
	\hspace{1pt}
	\psi'_{(\vec{\mathsf{x}}-\delta_m)/(\vec{\mathsf{y}}-\delta_m)}(t,q)
	\\
\hphantom{\Upsilon_m(\vec{\mathsf{x}},\vec{\mathsf{y}})=}{}	
\times
	\prod_{i=1}^{m}
	\big( 1-q^{m-i}t^{\mathsf{x}_i-m+i}\mathbf{1}_{\mathsf{y}_i=\mathsf{x}_i-1}\big)\nonumber
	\\
	\nonumber
\hphantom{\Upsilon_m(\vec{\mathsf{x}},\vec{\mathsf{y}})}{}	
=
	t^{-\binom m2}	
	\hspace{1pt}
	\prod_{1\le i<j\le m}
	\frac
	{
		\big(q^{j-i} t^{\mathsf{x}_i-\mathsf{x}_j-j+i+1};t\big)_{\infty}
		\big( q^{j-i-1}t^{\mathsf{y}_i-\mathsf{y}_j-j+i+1};t \big)_{\infty}
	}
	{
		\big( q^{j-i-1}t^{\mathsf{x}_i-\mathsf{x}_j-j+i+1};t \big)_{\infty}
		\big(q^{j-i} t^{\mathsf{y}_i-\mathsf{y}_j-j+i+1};t\big)_{\infty}
	}
	\\
\hphantom{\Upsilon_m(\vec{\mathsf{x}},\vec{\mathsf{y}})=}{}	\times
\nonumber
	\prod_{\substack{1\le i<j\le m
	\\
	\mathsf{y}_{i}=\mathsf{x}_{i}
	,\,
	\mathsf{y}_{j}=\mathsf{x}_{j}-1}}
	\frac{\big(1-t^{i-j+\mathsf{x}_{i}-\mathsf{x}_{j}+1}
	q^{j-i-1}\big)
	\big(1-t^{i-j+\mathsf{x}_{i}-\mathsf{x}_{j}} q^{j-i+1}\big)}
	{\big(1-t^{i-j+\mathsf{x}_{i}-\mathsf{x}_{j}+1} q^{j-i}\big)
	\big(1-t^{i-j+\mathsf{x}_{i}-\mathsf{x}_{j}} q^{j-i}\big)}
	\\
\hphantom{\Upsilon_m(\vec{\mathsf{x}},\vec{\mathsf{y}})=}{}	\times
	\prod_{i\colon \mathsf{y}_i=\mathsf{x}_i}
	t^{\mathsf{x}_i}
	\prod_{i\colon \mathsf{y}_i=\mathsf{x}_i-1}
	\big(t^{m-i}-q^{m-i}t^{\mathsf{x}_i}\big)
	.\label{eq:Upsilon_final_formula}
\end{gather}
The second expression is obtained by rewriting the
powers of $t$, and expanding the notation of~$V_{t,q}$ and~$\psi'$.
This completes the proof of Theorem~\ref{thm:Macdonald_noncolliding_intro}.
